\section{The Fisherman and the Keys}

No follower of Jesus of Nazareth is better documented than Simon Peter, and
none carries a heavier weight of later claim. The four Gospels name him more
often than every other disciple combined; Paul, writing within about
twenty-five years of the crucifixion, treats ``Cephas'' as a man his readers
already know; the Acts of the Apostles opens its story of the Church with his
voice. On the words Matthew places at Caesarea Philippi---``thou art Peter;
and upon this rock I will build my church\ldots\ And I will give to thee the
keys of the kingdom of heaven'' (Mt 16:18--19)---the Catholic Church rests
the doctrine of a permanent Petrine office, defined at the First Vatican
Council in 1870. Beneath the high altar of the largest church in Christendom,
excavators working under Pius XII uncovered a second-century memorial that a
Roman cleric named Gaius, around the year 200, had already called ``the
trophy'' of the apostle. A pope announced in 1968 that Peter's own bones had
been identified there ``in a manner we can consider convincing.''

Yet the same record is full of silences and reversals. The man the tradition
crowns is, in the earliest sources, the disciple who sank in the water,
misunderstood the Passion, was called ``Satan'' moments after being called
blessed, and denied three times that he knew Jesus at all. Scripture never
says that Peter went to Rome, never describes his death, and gives no date
for anything in his life. Everything after the Council of Jerusalem---the
Roman ministry, the martyrdom, the upside-down cross, the twenty-five-year
episcopate, the tomb---rests on witnesses of very unequal age and character,
from sober first-century allusion to frank second-century romance to
twentieth-century archaeology whose central identification is still
contested.

This biography therefore asks three questions. First, what does the earliest
controlled evidence---the Gospels read as distinct witnesses, Paul's letters,
and Acts---actually establish about Simon of Bethsaida, without a silent
harmony and without novelistic filling? Second, what can responsibly be said
about his end: the cumulative case for Rome and martyrdom under Nero, the
much later evidence for the manner of his death, and what the excavations
under St.~Peter's do and do not prove? Third, what have legend, liturgy, and
doctrine each made of him, and with what actual claims? The method is that of
the whole series: earliest evidence first; witnesses separated by decades or
centuries never silently merged; legends preserved and labeled rather than
suppressed or believed; doctrine reported as the Church's reception, not
converted into first-century data. The terminal appendix states the
conventions, chronology, and limits that govern every page.

Four questions in that program are genuinely contested, and this study does
not step around them. They are named here so that the reader knows in advance
where the argument will slow down and why.

\begin{itemize}
\item \textbf{The primacy texts.} Matthew 16:18--19, Luke 22:31--32, and
John 21:15--17 are the three passages on which the Petrine office is built.
They are read in their own contexts below, with the strongest non-Catholic
readings stated in the words of those who held them---Origen in the third
century, Calvin in the sixteenth---before any answer is attempted. The
exegesis is not settled by a biography; what a biography can do is show
exactly what the texts say, what they do not say, and which of the competing
readings the wording itself constrains.
\item \textbf{The Antioch incident.} Galatians 2:11--14 records the one
occasion when the first of the apostles was publicly rebuked by another for
conduct at variance with his own principle. Antiquity found the passage so
difficult that it produced three separate escape routes, and the last of
them---Jerome's---provoked a decade-long epistolary quarrel with Augustine
that survives in both men's letters. That quarrel is read here at the letters
themselves, because it is the earliest sustained Christian argument about
what apostolic authority does and does not guarantee.
\item \textbf{Rome and the martyrdom.} No first-century text places Peter in
Rome in so many words. The case is cumulative, and the honest way to test a
cumulative case is to state the case against it at full strength---the
silence of Paul's Letter to the Romans, the silence of Acts, the absence of
any narrative of the death---and then ask what the surviving evidence
actually carries.
\item \textbf{The letters.} 1~Peter and 2~Peter come to us under the
apostle's name, and their relation to the historical fisherman is not the
same in the two cases. They are graded separately below, with a statement of
which of this study's claims would fail if a given authorship judgment
turned out to be wrong.
\end{itemize}

Where this study reaches a judgment of its own rather than reporting a
source, it says so in a block headed \emph{Project synthesis}, states the
reasons, states the strongest argument against, and names the specific
finding that would change the judgment. Those blocks are the only places
where the study speaks in its own name about a disputed question.
